\subsection{Observable Leakage and Reachable Implementations}
\label{sec:leakage}
Target policy is one requirement; an exploitable observation is another.
A named weakness must
also be observable in the target's build and execution environment. The timing
(A1), cache (D1), binary-reachability (GOST), and micro-timing (E3) studies
examine this requirement. They establish the conditions under which a signal
or implementation is reachable; interpreting that signal as agent capability
requires additional evidence.
\subsubsection{Timing and Cache Observability}\label{sec:a1}
The A1 timing probe runs on two hosts (V tier, host-specific). On WSL2
(i7-9850H), 10{,}000 rounds yield a phase-2 mean of $-168$\,ns with
$z=1.91$; the pre-registered gate is FALSE. A false-positive control on the
same host yields $z=3.77$, showing that the gate correctly rejects a
non-separating distribution. On VMware (90{,}000 valid observations after
5\% trim of 100k per key), the mean is $-13.93$\,ns, median $-4$\,ns,
$\sigma=25.99$, SEM 0.1225, $z=113.70$, 95\% CI $[-14.17,-13.69]$, ASR 1,
TTE$_M$ 0.1\,s. The $z$ improvement is $59.5\times$ (tier E: derived from two
host statistics). Build, timer, preparation, and method change together
between hosts, so the contrast is confounded and exact key recovery remains
unresolved (Section~\ref{sec:limits}).
The E3 micro-timing factorial runs 16 cells $\times$ 2{,}000 trials
(32{,}000 trials, one bare-metal host, V tier) and yields an eight-cell paired
median of $-1.696$\,ms with 95\% CI $[-1.913,-1.590]$ and Welch
$p\approx8.9\times10^{-34}$. Warm medians are 0.761 vs.\ 2.350\,ms;
moving-block bootstrap replicates agree
($-1.721/-1.712/-1.696/-1.697$); a shakedown gate passes at $p=0.677$. A
15{,}971-record interrupted prefix is retained but excluded from the
eight-cell estimate. The factorial establishes a stable timing difference
under one host configuration; it does not transfer to the A1 hosts.
The D1 spread-table attack runs stages S0--S2 over 10 sessions, with 1/10
success (R tier; raw samples not retained). The spread-table arm collects
200{,}000 samples and recovers all 16 key bytes (XOR difference 0x00) against
a \emph{recompiled} victim that places each T-table entry on its own
4{,}096-byte page---a leakage-amplified fixture, not stock OpenSSL, on a
single Intel Skylake i7-6820HQ in user space. The Bernstein-style sibling
collects the same 200{,}000 samples in 281.45\,s and recovers 2/16 bytes at
rank 1 (8/16 in the top 10) against an ordinary-table victim. The two arms
compare two victim constructions within related cache-observation experiments;
Flush+Reload names a separate variant whose retained VM result is
inconclusive. The result is consistent with fixture layout affecting
observability, but the arms also differ in victim construction and execution
conditions, so it identifies neither one causal factor nor an agent effect.
\subsubsection{Implementation Reachability (Binary)}\label{sec:binary}
Two binary-reachability results separate recognizing a primitive from
executing it. The Windows \texttt{win64\_remote.exe} fallback (Runs \#3,
\#10) illustrates a reachable implementation path: the reported success
depends on the deployment state, while cryptanalytic derivation remains a
different requirement, and the retained summaries do not fully bind binary
digest + producing agent + generated exploit + outcome transcript. A GOST
investigation (P4-10) finds the string \texttt{GOST-28147} inside an
\emph{unreachable} legacy-compliance logging stub; recognizing a primitive
name does not establish execution, and this is a meaningful disconfirmation.
Legacy-crypto records distinguish four different things that earlier tables
conflated: static DES key extraction, weak-key dictionary, exhaustive-search
feasibility estimate ($>$1e16\,s, tier E), bounded Sweet32 collision, and RSA
factor verification. No unified ``legacy crypto'' rate is reported.